\begin{frame}
	\frametitle{Формулирование неопределённости}
	Рассмотрим динамику $\dot z^* = f(z^*, \zeta^*)$ с параметрами $\zeta^* = \text{const}$, и её разложение Тейлора в окрестности точки $z^* = z_0$, $\zeta^* = \zeta_0$. Обозначим $z = z^* - z_0$ и $\zeta = \zeta^* - \zeta_0$.
	Опустив квадратичные члены и члены высших порядков, но сохранив линейные и билинейные члены, запишем приближенную динамику в виде:
	
	\begin{equation}
		\label{eq:model_linearization_1}
		\dot z = \frac{\partial f}{\partial z} z + \Delta {A}(\zeta) z + \frac{\partial f}{\partial \zeta} \zeta,
	\end{equation}
	где $\Delta {A}(\zeta) = \frac{\partial}{\partial z} (\frac{\partial f}{\partial \zeta} \zeta)$.
	\vfill
	Модель \eqref{eq:model_linearization_1} включает в себя как мультипликативные, так и аддитивные неопределённости. Мы можем интерпретировать постоянные параметры $\zeta$ как неопределённые параметры модели, либо как статические компоненты состояния шагающего робота. 
\end{frame}

\begin{frame}
	\frametitle{Формулирование неопределённости}
	Предположим, что основным источником неопределённости является геометрия окружающей среды. Учитывая определение $\Delta {A}(\zeta)$, мы можем написать:
	\begin{equation}
		\Delta {A}(\zeta) = 
		\frac{\partial}{\partial z} 
		\left (
		\frac{\partial f}{\partial \zeta} {R}\T 
		\begin{bmatrix}
			{J}_K & 0 \\ \dot{{J}}_K & {J}_K
		\end{bmatrix}^+ 
		\begin{bmatrix}
			s_K \\ 0
		\end{bmatrix} \varsigma 
		\right )
		= \Tilde{A} \varsigma,
	\end{equation}
	%
	где  $-1 \leq \varsigma \leq 1$, $\Delta h$ --- неточность в положении точки на опорной поверхности, $h_K(q) = \Delta h$, ${J}_K = \partial h_K / \partial q$, $\Tilde{A}$ --- некоторая матрица, образующая базис в пространстве допустимых значений неопределённой матрицы $\Delta {A}$.
	
	Наша цель --- получить такие матрицы ${M}_1$, ${F}_1$ и ${N}_1$, где ${F}_1\T {F}_1 \leq {I}$, что для любого выбранного вектора $z$ множество всех векторов ${M}_1 {F}_1 {N}_1 z$ будет содержать множество $(\Tilde{A} \varsigma) z$. 
\end{frame}

\begin{frame}
    \frametitle{Математические инструменты}
\emph{S--процедура:}
Рассмотрим ${M},{N} \in \mathbb{R}^{n\times n}$. Если существует скаляр $\gamma>0$, который отвечает следующему условию ${M}+\gamma {N}<0$, тогда $x\T {N} x\geq 0$ подразумевает $x\T{M}x\leq 0$.

\emph{Неравенство Юнга:}
Для любой положительно определённой матрицы $M>0$, положительных скаляров $\epsilon > 0, \nu > 0$ и матриц ${X}, {Y}, {F}$, где ${F}\T{F}\leq \nu{I}$, следующие неравенства верны:
%
\begin{align}
	\label{eq:Young_relation_BMI}
	{X}\T{Y} + {Y}\T{X}  \leq {X}\T 
	\epsilon {M}^{-1}{X} + \frac{1}{\epsilon}   {Y}\T  {M}{Y}, 
	\\
	\label{eq:Young_robust_2}
	{X}\T{F}{Y} + {Y}\T{F}\T{X}  \leq \epsilon {X}\T{X} +  \frac{\nu}{\epsilon} {Y}\T {Y}.
\end{align}
%
Обозначим $\nu=\epsilon^2$ для формулирования следующего неравенства:
%
\begin{equation}
	\label{eq:updated_Young_robust_2}
	{X}^T{F}{Y} + {Y}\T{F}\T{X}  \leq \epsilon {X}\T{X} + \epsilon {Y}\T{Y}.
\end{equation}
\end{frame}

\begin{frame}
	\frametitle{Математические инструменты}
	\emph{Дополнение Шура:}
	Для любой симметричной матрицы ${A}\in \mathbb{S}^n$ и ${C}\in \mathbb{S}^m$ и матрицы ${B}\in \mathbb{R}^{n\times m}$, а также их конкатенации:
	\noindent \begin{align*}
		{M}= \begin{bmatrix}
			{A} & {B} \\
			{B}\T & {C} 
		\end{bmatrix},
	\end{align*}
	%
	следующие утверждения равнозначны:
	% 
	\noindent
	\begin{enumerate}
		\item ${M} < 0$,
		\item ${A}-{B}{C}^{-1}{B}\T < 0 , {C}< 0$,
		\item ${C}-{B}\T{A}^{-1}{B}< 0 , {A}< 0$.
	\end{enumerate}
\end{frame}

\begin{frame}
	\frametitle{Целевая функция}
	Неопределённости ограничены следующим образом:
	%
	\begin{equation}
		{F}_i\T{F}_i\leq \nu_i {I}, \ \ i=1,2,3,
	\end{equation}
	где $\nu_i$ --- радиусы неопределённости --- скаляры, определяющие ограничения, накладываемые на нормы ${F}_i$. 
	
	Обозначим $\mu_i=\gamma_i\nu_i$ и $\bar{\mu}_i=\frac{1}{\mu_i}$ для $i=1,2,3$.
	Используя данные переменные, можно сформулировать следующие целевые функции:
	\begin{align}
		\label{eq:cost_lin}
		J_\text{lin} &= \sum_{i=1}^{3}\left(\bar{\mu}_i+\gamma_i\right), \\ 
		\label{eq:cost_quad}
		J_\text{quad} & =  \sum_{i=1}^{3}\left((\bar{\mu}_i-\varpi)^2+(\gamma_i-\varpi)^2\right) + \varpi^2,
	\end{align}
	%
	где $\varpi$ --- свободная переменная. Целью обеих целевых функций является максимизация $\nu_i$, которая достигается минимизацией либо $J_\text{lin}$, либо $J_\text{quad}$. Максимизируя $\nu_i$, достигается устойчивость к большему набору неопределённых матриц $\Delta {A}_n$, $\Delta {B}$ и $\Delta {C}$.
\end{frame}

\begin{frame}
	\frametitle{Системы с механическими ограничениями}
    \begin{figure}
	\begin{minipage}[c]{0.4\textwidth}
		\textbf{Ортогональные методы}
		\begin{equation*}
			\begin{gathered}
				\dot{x}=Ax+Bu+F\lambda\\
				G\dot{x}=0
			\end{gathered}
		\end{equation*}
	\end{minipage}\hfill
	\begin{minipage}[c]{0.6\textwidth}
		\textbf{Дескрипторные методы}
		\begin{equation*}
			\begin{gathered}
				E\dot{x}=Ax+Bu\\
				\begin{pmatrix}
					I & 0\\
					G & 0
				\end{pmatrix}
				\begin{pmatrix}
					\dot{x}\\ \dot{\lambda}
				\end{pmatrix}=
				\begin{pmatrix}
					A & F\\ 0 & 0
				\end{pmatrix}
				\begin{pmatrix}
					x\\ \lambda
				\end{pmatrix}+
				\begin{pmatrix}
					B \\ 0
				\end{pmatrix}
				u
			\end{gathered}
		\end{equation*}
	\end{minipage}
\end{figure}
\end{frame}

\begin{frame}
	\frametitle{Применение методов}
	\begin{figure}
		\centering
		\includegraphics[width=0.8\linewidth]{images/output_state_z.JPG}
		\includegraphics[width=0.8\linewidth]{images/output_state_zk.JPG}
		\caption{Вектор состояния и вектор состояния регулятора от времени при решении \eqref{eq:thm2_OCP} для плоского четвероногого робота.}
		\label{fig:output_state_z}
	\end{figure}
\end{frame}

\begin{frame}
	\frametitle{Предмет и объект исследования}
	Объектом исследования является оптимальное робастное управление шагающими роботами с параметрическими неопределённостями.
	\vfill
	Предметом исследования являются методы математического моделирования системы с неопределённостями, описывающие динамику шагающего робота и численные методы с использованием ортогональных методов и линейных матричных неравенств для нахождения коэффициентов регулятора и наблюдателя с последующими вычислениями при помощи комплекса программ.
\end{frame}